\paragrafos{Τυποποιημένη μορφή: οι μεγάλοι αριθμοί χωράνε σε μια γραμμή}{Α.3.4, «Τυποποιημένη μορφή μεγάλων αριθμών»}

\begin{anagke}
Η κυρία Αλεξάνδρα ζητάει από την Α3 μια εργασία για το ηλιακό σύστημα. Η Αγγελική γράφει στο τετράδιό της την απόσταση της Γης από τον Ήλιο: $149\,600\,000$ χιλιόμετρα. Ο Ηλίας γράφει τη μάζα της Γης: $5\,976\,000\,000\,000\,000\,000\,000\,000$ κιλά, και το γράφει τρεις φορές, γιατί δεν είναι σίγουρος για τα μηδενικά. Ο Στέλιος δοκιμάζει με τον Ήλιο και χάνεται στο δεύτερο μέτρημα. «Αυτό είναι αδύνατο να διαβαστεί», παραπονιέται η Μυρτώ.

Η Πολυξένη, που έχει το κομπιουτεράκι της, πατάει $149\,600\,000\cdot10$ και η οθόνη δείχνει $1.496\mathrm{E}9$. «Τι είναι αυτό το E9;». Η κυρία Αλεξάνδρα χαμογελάει: «Το κομπιουτεράκι δεν χωράει όλα τα μηδενικά, και τα γράφει με συντομία: $1{,}496$ επί δέκα στην ένατη. Κάνουν ό,τι κάνουν και οι επιστήμονες». Και η Φανή ρωτάει: «Μπορούμε να γράφουμε κι εμείς έτσι, ώστε να μη χάνουμε τα μηδενικά;».
\end{anagke}

\begin{lexi}{τυποποιημένη μορφή}{τυποποιημένη μορφή}
Η λέξη \lexiword{τύπος} σημαίνει καλούπι, πρότυπο. \lexiword{Τυποποιημένη} είναι η μορφή που έχει το ίδιο καλούπι για όλους τους αριθμούς: ένας αριθμός από $1$ μέχρι $9{,}99\ldots$ επί μια δύναμη του $10$. Οι επιστήμονες την χρησιμοποιούν γιατί δείχνει αμέσως το \emph{μέγεθος} του αριθμού (πόσα μηδενικά) και τα σημαντικά του ψηφία.
\end{lexi}

\begin{metafrasi}
Εκατόν σαράντα εννέα εκατομμύρια εξακόσιες χιλιάδες. & $149\,600\,000=1{,}496\cdot10^8$\\
Τόσα μηδενικά όσο ο εκθέτης. & $10^8=100\,000\,000$ (το $1$ και $8$ μηδενικά)\\
Ο $4{,}56\cdot10^{10}$ γραμμένος κανονικά. & $45\,600\,000\,000$\\
Στην οθόνη του κομπιουτεράκι. & $1.496\mathrm{E}8$ σημαίνει $1{,}496\cdot10^8$\\
Σύγκριση: μεγαλύτερος εκθέτης. & $8{,}2\cdot10^5>9{,}1\cdot10^4$\\
\end{metafrasi}

\tmima{Η θεωρία}

\begin{orismos}{τυποποιημένη μορφή}
Ένας (μεγάλος) αριθμός είναι γραμμένος σε \textbf{τυποποιημένη μορφή} όταν είναι γραμμένος ως $\alpha\cdot10^{\nu}$, όπου $\nu$ φυσικός αριθμός και ο $\alpha$ είναι ένας δεκαδικός (ή φυσικός) με ακέραιο μέρος ένα ψηφίο, δηλαδή $1\leq\alpha<10$.
\end{orismos}

Για παράδειγμα
$$3\,140\,000\,000\,000\,000\,000=3{,}14\cdot10^{18},\qquad 234\,000\,000\,000\,000\,000=2{,}34\cdot10^{17}.$$
Και αντίστροφα, ο $5{,}21\cdot10^5$ είναι ο $521\,000$ και ο $2\cdot10^3$ ο $2\,000$.

\tmimasmall{Πώς γράφω έναν αριθμό σε τυποποιημένη μορφή}

\begin{kanonas}{τυποποίηση}
\begin{enumerate}[leftmargin=1.6em,itemsep=0.3ex,topsep=2pt]
\item Μετακινώ την υποδιαστολή (στους ακεραίους είναι στο τέλος) προς τα αριστερά, μέχρι να μείνει \textbf{ένα μη μηδενικό ψηφίο} πριν από αυτήν. Έτσι βρίσκω τον $\alpha$ (με $1\leq\alpha<10$).
\item Ο εκθέτης $\nu$ είναι το πλήθος των θέσεων που μετακίνησα.
\item Γράφω $\alpha\cdot10^{\nu}$. Αφαιρώ τα μηδενικά που περισσεύουν στο τέλος του $\alpha$.
\end{enumerate}
\end{kanonas}

Για να πάω πίσω, από την τυποποιημένη στη δεκαδική μορφή, πολλαπλασιάζω τον $\alpha$ επί $10^{\nu}$: μετακινώ την υποδιαστολή του $\alpha$ δεξιά κατά $\nu$ θέσεις, βάζοντας μηδενικά όπου χρειάζεται. Έτσι $7{,}4\cdot10^3=7\,400$.

\tmimasmall{Σύγκριση και πράξεις}

Για να συγκρίνω δύο αριθμούς σε τυποποιημένη μορφή, κοιτάζω πρώτα τον \textbf{εκθέτη}: μεγαλύτερος εκθέτης σημαίνει μεγαλύτερο αριθμό, επειδή έχει περισσότερα ψηφία. Αν οι εκθέτες είναι ίσοι, συγκρίνω τους $\alpha$.

Για να πολλαπλασιάσω, πολλαπλασιάζω τους $\alpha$ μεταξύ τους και προσθέτω τους εκθέτες του $10$: $(2{,}5\cdot10^4)\cdot(4\cdot10^3)=10\cdot10^7=1\cdot10^8$. Προσοχή: αν το γινόμενο των $\alpha$ βγει $10$ ή περισσότερο, πρέπει να ξαναγράψω το αποτέλεσμα σε τυποποιημένη μορφή.

\begin{prosoxi}
\begin{itemize}[leftmargin=1.2em,itemsep=0.3ex,topsep=2pt]
\item Ο $\alpha$ είναι πάντα μικρότερος από το $10$ και τουλάχιστον $1$. Οι $52\cdot10^2$, $0{,}52\cdot10^4$ και $10\cdot10^5$ δεν είναι τυποποιημένοι (αν και ισούνται με $5{,}2\cdot10^3$, $5{,}2\cdot10^3$ και $1\cdot10^6$ αντίστοιχα).
\item Ο εκθέτης δεν είναι το πλήθος των ψηφίων αλλά το πλήθος των ψηφίων μείον ένα. Ο $45\,600\,000\,000$ έχει $11$ ψηφία και γράφεται $4{,}56\cdot10^{10}$.
\item Δεν συγκρίνω τους $\alpha$ όταν διαφέρουν οι εκθέτες: $9{,}9\cdot10^3<1\cdot10^4$, αν και το $9{,}9>1$.
\end{itemize}
\end{prosoxi}

\tmima{Λυμένα παραδείγματα}

\paradeigma{Η απόσταση Γη -- Ήλιος}
\tixero Η Γη απέχει από τον Ήλιο $149\,600\,000$ χλμ.\par
\tithelo Να το γράψω σε τυποποιημένη μορφή.\par
\possindeo Μετακινώ την υποδιαστολή (που είναι στο τέλος) προς τα αριστερά ώσπου να μείνει ένα ψηφίο πριν από αυτήν: $1{,}49600000$. Μετακίνησα $8$ θέσεις. Αφαιρώ τα τελευταία μηδενικά.\par
\apantisi $1{,}496\cdot10^8$ χλμ.\par
\elegxos $1{,}496\cdot100\,000\,000=149\,600\,000$. Και ο αριθμός έχει $9$ ψηφία, εκθέτης $9-1=8$.

\paradeigma{Από τυποποιημένη σε δεκαδική}
\tixero Οι αριθμοί $7{,}03\cdot10^6$ και $4{,}56\cdot10^{10}$.\par
\tithelo Να τους γράψω κανονικά.\par
\possindeo Ο $7{,}03\cdot10^6$: μετακινώ την υποδιαστολή $6$ θέσεις δεξιά, $7\,030\,000$. Ο $4{,}56\cdot10^{10}$: $10$ θέσεις δεξιά, $45\,600\,000\,000$.\par
\apantisi $7\,030\,000$ και $45\,600\,000\,000$.\par
\elegxos Στον πρώτο υπάρχουν $7$ ψηφία (εκθέτης $6$) και στον δεύτερο $11$ ψηφία (εκθέτης $10$).

\paradeigma{Ποιος είναι μεγαλύτερος;}
\tixero Οι αριθμοί $8{,}2\cdot10^5$, $9{,}1\cdot10^4$, $3{,}5\cdot10^7$ και $2{,}9\cdot10^7$.\par
\tithelo Να τους βάλω σε αύξουσα σειρά.\par
\possindeo Κοιτάζω πρώτα τους εκθέτες: $4<5<7$. Ο $9{,}1\cdot10^4$ είναι ο μικρότερος και ο $8{,}2\cdot10^5$ ακολουθεί. Οι άλλοι δύο έχουν εκθέτη $7$, και $2{,}9<3{,}5$.\par
\apantisi $9{,}1\cdot10^4<8{,}2\cdot10^5<2{,}9\cdot10^7<3{,}5\cdot10^7$.\par
\elegxos Δεκαδικά: $91\,000<820\,000<29\,000\,000<35\,000\,000$.

\paradeigma{Πόση ώρα ταξιδεύει το φως;}
\tixero Το φως ταξιδεύει με $3\cdot10^5$ χλμ. ανά δευτερόλεπτο και ο Ήλιος απέχει $1{,}496\cdot10^8$ χλμ.\par
\tithelo Σε πόσο χρόνο το φως φτάνει από τον Ήλιο στη Γη.\par
\possindeo Χρόνος = απόσταση : ταχύτητα $=\dfrac{1{,}496\cdot10^8}{3\cdot10^5}=\dfrac{1{,}496}{3}\cdot10^3\approx0{,}4987\cdot10^3\approx499$ δευτερόλεπτα. Στα λεπτά: $499:60\approx8{,}3$.\par
\apantisi Περίπου $8{,}3$ λεπτά.\par
\elegxos Σε $8{,}3$ λεπτά, δηλαδή περίπου $499$ δευτερόλεπτα, το φως διανύει $499\cdot300\,000\approx1{,}5\cdot10^8$ χλμ. Συμφωνεί.

\paradeigma{Γινόμενο σε τυποποιημένη μορφή}
\tixero Οι αριθμοί $2{,}5\cdot10^4$ και $4\cdot10^3$.\par
\tithelo Το γινόμενό τους σε τυποποιημένη μορφή.\par
\possindeo Πολλαπλασιάζω τους $\alpha$: $2{,}5\cdot4=10$. Προσθέτω τους εκθέτες: $4+3=7$. Έχω $10\cdot10^7$. Το $10$ δεν είναι μικρότερο του $10$, γι' αυτό γράφω $10\cdot10^7=10^8$.\par
\apantisi $1\cdot10^8$.\par
\elegxos $25\,000\cdot4\,000=100\,000\,000=10^8$.

\tmima{Ασκήσεις}

\zoni{Το κατάλαβα}{τυποποίηση, αποτυποποίηση, σύγκριση}

\ex Γράψε σε τυποποιημένη μορφή: α) $72\,000$ \quad β) $1\,250\,000$ \quad γ) $903\,000\,000$ \quad δ) $6\,000\,000\,000$ \quad ε) $48\,700$ \quad στ) $15\,000\,000\,000\,000$.

\ex Γράψε με δεκαδική μορφή: α) $2{,}7\cdot10^5$ \quad β) $5{,}06\cdot10^3$ \quad γ) $1{,}9\cdot10^7$ \quad δ) $8\cdot10^4$ \quad ε) $3{,}125\cdot10^6$ \quad στ) $4{,}008\cdot10^9$.

\ex Ποιοι από τους αριθμούς είναι σε τυποποιημένη μορφή; α) $5{,}2\cdot10^3$ \quad β) $52\cdot10^2$ \quad γ) $0{,}52\cdot10^4$ \quad δ) $10\cdot10^5$ \quad ε) $9{,}99\cdot10^1$.

\ex Γράψε σε τυποποιημένη μορφή: α) $24\cdot10^5$ \quad β) $0{,}35\cdot10^7$ \quad γ) $150\cdot10^3$ \quad δ) $0{,}008\cdot10^9$.

\ex Σύγκρινε με τα $<$, $=$, $>$: α) $8{,}2\cdot10^5$ και $9{,}1\cdot10^4$ \quad β) $3{,}5\cdot10^7$ και $2{,}9\cdot10^7$ \quad γ) $6\cdot10^9$ και $5{,}99\cdot10^9$ \quad δ) $4{,}4\cdot10^3$ και $4{,}4\cdot10^5$.

\ex Γράψε σε αύξουσα σειρά: $3{,}2\cdot10^6$, $9\cdot10^5$, $1{,}1\cdot10^7$, $4{,}5\cdot10^6$, $7{,}9\cdot10^5$.

\ex Η Πολυξένη βλέπει στην οθόνη του κομπιουτεράκι: α) $4.5\mathrm{E}7$ \quad β) $2.5\mathrm{E}10$. Ποιοι είναι οι αριθμοί; Πόσο δείχνει η οθόνη αν πολλαπλασιάσει τον $12\,000\,000$ με τον $3\,000\,000$;

\zoni{Το χρησιμοποιώ}{ηλιακό σύστημα και ανθρώπινο σώμα}

\ex \textbf{Σελήνη και Ήλιος.} Η Σελήνη απέχει $384\,400$ χλμ. από τη Γη και ο Ήλιος $149\,600\,000$ χλμ. Γράψε τις αποστάσεις σε τυποποιημένη μορφή και πες περίπου πόσες φορές μακρύτερα είναι ο Ήλιος.

\ex \textbf{Το έτος φωτός.} Ένα έτος φωτός είναι $9\,461\,000\,000\,000$ χλμ. Γράψε το σε τυποποιημένη μορφή. Ο κοντινότερος αστέρας, ο Εγγύτατος του Κενταύρου, απέχει περίπου $4{,}3$ έτη φωτός. Πόσα χιλιόμετρα είναι αυτά, σε τυποποιημένη μορφή, με δύο δεκαδικά;

\ex \textbf{Το φως του Ήλιου.} Η Φανή ρωτάει σε πόσο χρόνο φτάνει στη Γη το φως, αν ταξιδεύει με $3\cdot10^5$ χλμ. το δευτερόλεπτο και ο Ήλιος απέχει $1{,}496\cdot10^8$ χλμ.; Απάντησε σε δευτερόλεπτα και σε λεπτά.

\ex \textbf{Πληθυσμοί.} Ο πληθυσμός της Γης είναι περίπου $8\,000\,000\,000$ άνθρωποι και της Ελλάδας περίπου $10\,400\,000$. Γράψε τους σε τυποποιημένη μορφή και βρες περίπου πόσες φορές μεγαλύτερος είναι ο πρώτος.

\ex \textbf{Οι μάζες.} Η μάζα της Γης είναι $5{,}976\cdot10^{24}$ κιλά και του Ήλιου $1{,}983\cdot10^{30}$ κιλά. Γράψε τη μάζα της Γης σε δεκαδική μορφή και βρες περίπου πόσες φορές είναι ο Ήλιος βαρύτερος.

\ex \textbf{Η μνήμη του κινητού.} Η μνήμη του κινητού της Μυρτώς είναι $128$ GB, δηλαδή $1{,}28\cdot10^{11}$ bytes. Μια φωτογραφία έχει περίπου $4$ MB, δηλαδή $4\cdot10^6$ bytes. Πόσες φωτογραφίες χωράνε;

\zoni{Το σκέφτομαι}{οι κανόνες, όταν σπάνε}

\ex Γιατί ο $10\cdot10^5$ δεν είναι σε τυποποιημένη μορφή; Πώς γράφεται σωστά;

\ex Πόσα ψηφία έχει ο $4{,}56\cdot10^{10}$ όταν γραφεί κανονικά; Και ο $1\cdot10^{\nu}$; Εξήγησε πώς το βρίσκεις από τον εκθέτη.

\ex Υπολόγισε τα $(3\cdot10^4)\cdot(2\cdot10^5)$ και $(4\cdot10^3)\cdot(5\cdot10^6)$. Βάλε τα αποτελέσματα σε τυποποιημένη μορφή και διατύπωσε τον κανόνα του πολλαπλασιασμού.

\ex Ο Βασίλης λέει ότι $9{,}9\cdot10^3>1\cdot10^4$, γιατί $9{,}9>1$. Έχει δίκιο;

\ex Αν μετράμε αριθμούς φωναχτά, ένα νούμερο ανά δευτερόλεπτο, πόσα χρόνια χρειαζόμαστε για να φτάσουμε στο $10^{15}$; (Ένα έτος έχει περίπου $3{,}15\cdot10^7$ δευτερόλεπτα.)

\zoni{Για το σπίτι}{ψάξε τους μεγάλους αριθμούς του κόσμου}

\ex Βρες από αξιόπιστη πηγή τη μάζα του Ηλίου, το πλήθος των αστεριών του γαλαξία μας και την ακτίνα της Γης. Γράψε τους σε δεκαδική και σε τυποποιημένη μορφή.

\ex Φτιάξε πέντε δικούς σου αριθμούς, γράψε τους σε τυποποιημένη μορφή και ζήτησε από έναν συμμαθητή να βρει τους αρχικούς. Έλεγξε τις απαντήσεις του.
